\begin{abstract}
Language models are adept at conversation, offering advice and suggestions
with apparent empathy and understanding. Whatever the nature of that
understanding, it is widely accepted to break down under more complex social
conditions, and the sharpest case is sycophancy, the answer pulled toward the
person asking. The social account of objectivity holds that organised
criticism across perspectives exposes errors that no single mind corrects,
and on that account we introduce a general, lightweight prompt scaffold
structured as a narrative. Narration-of-Thought (NoT) guides the model
through a first-person narrative that identifies the decision-maker and what
they know, considers every affected party and their interests, traces the
consequences of the available actions, acknowledges uncertainty, and only
then commits. We show that the scaffold can be used in single-agent and
multi-agent settings on a range of nuanced social reasoning tasks. Used alone,
it reduces validation of the asker's own account on open-ended personal
dilemmas by 22 to 70 points
relative to standard chain-of-thought on seven models from six vendors,
three of them open-weight, a relative reduction of up to 87 percent,
and the reduction is three times as large on accounts that name a second
party.
\todo{Pillar two (theory of mind) is a skeleton for a future developer and
is not claimed in this paper; no sentence on it belongs in the abstract
until it is in scope.}
Used to build a collective, copies of one model given opposing sides
support more objective, grounded judgements through structured criticism,
and reviewing the cases flagged by a second dissenting seat with a stronger
judge that has no side raises accuracy by 7.0 points over the collective
alone
and by 2.4 points over the strongest single model
in our primary evaluation. We map where each effect holds and where it does
not, and our findings lay the groundwork for further work on the scaffold as
an instrument of social alignment.
\end{abstract}
